% !TeX root = ../../Book4.tex
% 纲要标题：### 21.1 DG 代数与 DG 范畴
% 纲要标签：b4:sec:2001
% 纲要位置：计划纲要.md，第 3026 行。

\section{DG 代数与 DG 范畴}
\label{b4:sec:2001}

在适当的投射或内射模型上，导出态射由复形映射的同伦类表示。完整 Hom 复形还保留具体同伦及其复合；将这些运算放在一起，乘法与微分的符号便须满足分次 Leibniz 法则。


% 纲要内容（仅作写作计划，尚非教材正文）：
%
% * 区别第三卷附录 E 的指定导子：这里采用分次微分、带符号 Leibniz 法则及微分平方为零
% * 带微分的分次态射空间
% * 复合与符号
% * 同伦范畴及增强
%

\subsection{分次微分与带符号 Leibniz 法则}
\label{b4:subsec:2001:01}

先看一个可以完全算出的 Hom 复形。设 \(k\) 为域，\(C=(k\xrightarrow{1}k)\) 位于零次和一次。它可缩，所以在同伦范畴中为零对象。不过，次数 \(-1\) 的映射 \(h:C^1\to C^0\)、\(h=1\)，仍记录了收缩同伦。其分次自同态空间 \(E=\underline{\operatorname{End}}_k(C)\) 连同微分为
\[
 E^{-1}=k\xrightarrow{\ c\mapsto(c,c)\ }E^0=k\oplus k
 \xrightarrow{\ (a,b)\mapsto a-b\ }E^1=k.
\]
这里微分由 \(\partial f=\mathrm{d}_Cf-(-1)^{|f|}f\mathrm{d}_C\) 给出，尤其 \(\partial h=1_C\)。三处上同调都为零，但 \(E\) 还保留了使恒等映射零伦的具体 \(h\)，以及所有映射的复合。要同时保存这些数据，就须把分次乘法与微分放在同一结构中。

\begin{definition}\label{b4:def:dg-algebra}
设 \(k\) 为交换环，\(A=\bigoplus_{n\in\mathbb Z}A^n\) 为含幺的分次结合 \(k\)-代数，单位位于 \(A^0\)。设 \(\mathrm{d}:A^n\to A^{n+1}\) 为 \(k\)-线性映射，满足 \(\mathrm{d}^2=0\)，并且对 \(a\in A^p,b\in A^q\) 有
\[
 \mathrm{d}(ab)=\mathrm{d}(a)b+(-1)^p a\,\mathrm{d}(b).
\]
称 \((A,\mathrm{d})\) 为\term[weifenfencidaishu]{微分分次代数}[differential graded algebra]，简称 DG 代数。其态射保持分次、乘法、单位与微分。
\end{definition}

由 \(\mathrm{d}(1)=\mathrm{d}(1\cdot1)=2\mathrm{d}(1)\) 得 \(\mathrm{d}(1)=0\)，这里不需要除以 \(2\)。符号取决于左因子的次数：微分越过一个 \(p\) 次元素时出现 \((-1)^p\)。普通 \(k\)-代数放在零次、微分取零，便成为 DG 代数；DG 代数则不要求乘法分次交换。

\ProseParagraphBreak
第三卷附录 E 的微分代数从一个满足普通 Leibniz 法则的指定导子出发，导子不一定平方为零。本节的微分升高次数，既有符号条件，又必须平方为零。例如 \(k[x]\) 上的 \(\partial/\partial x\) 并不因名称中含有“微分”就给出这里的结构；在特征零时，它作用于 \(x^2\) 两次的结果为 \(2\)。

\begin{example}\label{b4:exm:dg-koszul-one-element}
设 \(R\) 为交换环，\(a\in R\)。令 \(A^0=R\)、\(A^{-1}=R\varepsilon\)，其余次数为零，规定 \(\varepsilon^2=0\)、\(\mathrm{d}\varepsilon=a\)、\(\mathrm{d}R=0\)。这给出一个 DG \(R\)-代数。唯一需要特别检查的关系是
\[
 \mathrm{d}(\varepsilon^2)=a\varepsilon-\varepsilon a=0.
\]
其上同调为 \(H^0(A)=R/aR\)、\(H^{-1}(A)=\operatorname{ann}_R(a)\)。当 \(a\) 不是零因子时，增广 \(A\to R/aR\) 是拟同构；当 \(a\) 为单位时，底层复形可缩。这里重现了单元素 Koszul 复形，但同时记录了外代数乘法。
\end{example}

\begin{proposition}\label{b4:prop:dg-cohomology-product}
设 \(A\) 为交换环 \(k\) 上的 DG 代数。公式 \([a][b]=[ab]\) 在 \(H^*(A)=\bigoplus_n H^n(A)\) 上给出分次结合代数结构。
\end{proposition}
\begin{proof}
若 \(\mathrm{d}a=\mathrm{d}b=0\)，则 Leibniz 法则给出 \(\mathrm{d}(ab)=0\)。若把 \(a\) 改为 \(a+\mathrm{d}u\)，其中 \(u\) 为 \(p-1\) 次元素，则乘积之差为 \((\mathrm{d}u)b=\mathrm{d}(ub)\)；若把 \(b\) 改为 \(b+\mathrm{d}v\)，则 \(a(\mathrm{d}v)=(-1)^p \mathrm{d}(av)\)。因此乘积不依赖余圈代表。结合律和单位由 \(A\) 中的对应等式下降；若 \([1]=0\)，所得上同调代数为零代数。
\end{proof}

\subsection{带微分的分次态射空间}
\label{b4:subsec:2001:02}

设 \(R\) 为含幺环，\(C,D\) 为左 \(R\)-模上链复形。一个次数为 \(r\) 的齐次映射，是一族不必与微分交换的 \(R\)-线性映射 \(f^n:C^n\to D^{n+r}\)。全部这类映射组成
\[
 \underline{\operatorname{Hom}}_R(C,D)^r
 =\prod_{n\in\mathbb Z}\operatorname{Hom}_R(C^n,D^{n+r}).
\]
积号不能随意改为直和：一个齐次映射可以在无限多个次数上非零。令
\begin{equation}\label{b4:eq:dg-hom-differential}
 (\partial f)^n=\mathrm{d}_D^{n+r}f^n-(-1)^r f^{n+1}\mathrm{d}_C^n.
\end{equation}
这正是前面总 Hom 复形的微分；此处把同一公式用于组织所有复形之间的映射。

\begin{proposition}\label{b4:prop:dg-hom-complex}
设 \(R\) 为含幺环，\(C,D\) 为左 \(R\)-模上链复形。在 \(\underline{\operatorname{Hom}}_R(C,D)\) 上，对次数 \(r\) 的齐次映射 \(f\) 定义 \(\partial f=\mathrm d_Df-(-1)^rf\mathrm d_C\)，则 \(\partial^2=0\)。其零次余圈是复形映射，零次余边是零伦映射，因而
\[
 H^0\underline{\operatorname{Hom}}_R(C,D)
 \cong\operatorname{Hom}_{K(R\text{-}\mathrm{Mod})}(C,D).
\]
\end{proposition}
\begin{proof}
省略上标而保留次数，展开两次微分得到
\[
 \partial^2 f=\mathrm{d}_D^2f-(-1)^r \mathrm{d}_Df\mathrm{d}_C
             -(-1)^{r+1}\mathrm{d}_Df\mathrm{d}_C-f\mathrm{d}_C^2=0.
\]
在零次，\(\partial f=0\) 就是 \(\mathrm{d}_Df=f\mathrm{d}_C\)。若 \(h\) 的次数为 \(-1\)，则 \(\partial h=\mathrm{d}_Dh+h\mathrm{d}_C\)，恰是零伦公式。取商即得到所述同伦范畴中的态射群。
\end{proof}

较高次数也有明确意义：\(\partial f=0\) 等价于把 \(f^n\) 看作 \(C^n\to D[r]^n\) 后得到复形映射，因为 \(\mathrm{d}_{D[r]}=(-1)^r \mathrm{d}_D\)。相应余边与同伦只差一个整体符号，故 \(H^r\underline{\operatorname{Hom}}_R(C,D)=\operatorname{Hom}_{K(R\text{-}\mathrm{Mod})}(C,D[r])\)。若 \(C\) 是有界投射复形，\cref{b4:thm:bounded-projective-model}又把这里的同伦态射识别为导出态射。任意复形则不能省去这个条件。

\subsection{复合与符号}
\label{b4:subsec:2001:03}

齐次映射仍按通常次序复合。若 \(f:C\to D\) 的次数为 \(p\)，\(g:D\to E\) 的次数为 \(q\)，则 \((gf)^n=g^{n+p}f^n\)，次数为 \(p+q\)。直接展开可得
\begin{equation}\label{b4:eq:dg-composition-leibniz}
 \partial(gf)=(\partial g)f+(-1)^q g(\partial f).
\end{equation}
右边的中间两项为 \(-(-1)^q g\mathrm{d}_Df\) 与 \(( -1)^q g\mathrm{d}_Df\)，正好抵消。因此符号由写在左边、后执行的 \(g\) 的次数决定，而不是由先执行的 \(f\) 决定。

\begin{definition}\label{b4:def:dg-category}
设 \(k\) 为交换环。给定一组数据 \(\mathcal E\)：一族对象；对每两个对象 \(X,Y\)，给定 \(k\)-模上链复形 \(\underline{\operatorname{Hom}}_{\mathcal E}(X,Y)\)，称其为态射复形；给定双线性、保持总次数的复合，以及零次余圈 \(1_X\in\underline{\operatorname{Hom}}_{\mathcal E}(X,X)^0\)。若复合结合，\(1_X\) 是单位，并且对次数 \(p,q\) 的齐次态射 \(f:X\to Y\)、\(g:Y\to Z\) 满足 \(\mathrm d(gf)=(\mathrm dg)f+(-1)^qg(\mathrm df)\)，则称这些数据构成一个\term[weifenfencifanchou]{微分分次范畴}[differential graded category]，简称 DG 范畴。

两个 \(k\)-线性 DG 范畴之间，若对象对应与各 Hom 复形之间的复形映射保持单位和复合，称其为\term[DGhanzi]{DG 函子}[DG functor]。
\end{definition}

\(R\)-模复形及上述 Hom 复形组成 DG 范畴 \(\mathrm{Ch}_{\mathrm{dg}}(R)\)。只有一个对象的 DG 范畴等价于一个 DG 代数；特别是 \(\underline{\operatorname{End}}_R(C)=\underline{\operatorname{Hom}}_R(C,C)\) 以复合作为乘法，是 DG 代数。其零次上同调才是 \(C\) 的同伦自同态环，二者一般不是同一个对象。

\ProseParagraphBreak
张量积中出现的符号有相同来源。设 \(A,B\) 为交换环 \(k\) 上的 DG 代数，对齐次元素规定
\[
 (a\otimes b)(a'\otimes b')=(-1)^{|b|\cdot|a'|}aa'\otimes bb',
 \qquad
 \mathrm{d}(a\otimes b)=\mathrm{d}a\otimes b+(-1)^{|a|}a\otimes \mathrm{d}b.
\]
在展开乘积的微分时，把 \(b\) 越过 \(a'\) 所需的符号恰好使 Leibniz 法则成立。与此相配，DG 反代数的乘法为 \(a\mathbin{\cdot_{\mathrm{op}}}b=(-1)^{|a|\cdot|b|}ba\)，而不只是反转次序。这些约定在 DG 双模和张量函子中不可省略；集中在零次时，它们还原为普通约定。

\subsection{同伦范畴及增强}
\label{b4:subsec:2001:04}

对 DG 范畴 \(\mathcal E\)，保留对象而把态射复形换成零次上同调，得到范畴 \(H^0(\mathcal E)\)，其中
\[
 \operatorname{Hom}_{H^0(\mathcal E)}(X,Y)
 =H^0\underline{\operatorname{Hom}}_{\mathcal E}(X,Y).
\]
复合良定义，是因为余圈复合仍是余圈，而且一个零次余边与零次余圈复合后仍是余边。\cref{b4:prop:dg-hom-complex}说明
\[
 H^0\bigl(\mathrm{Ch}_{\mathrm{dg}}(R)\bigr)=K(R\text{-}\mathrm{Mod}).
\]
因此取 \(H^0\) 已经忘掉了同伦的具体选择。它还没有自动把所有拟同构倒置。

\begin{definition}\label{b4:def:dg-quasi-equivalence}
设 \(F:\mathcal E\to\mathcal F\) 为交换环 \(k\) 上的 DG 函子。若每个 \(\underline{\operatorname{Hom}}_{\mathcal E}(X,Y)\to\underline{\operatorname{Hom}}_{\mathcal F}(FX,FY)\) 是拟同构，而且 \(H^0(F)\) 本质满，称 \(F\) 为\term[nidengjia]{拟等价}[quasi-equivalence]。
\end{definition}

拟等价诱导 \(H^0\) 上的等价，因为零次上同调上的映射已经双射。反过来，仅知道零次上同调范畴等价，并没有提供高次映射或与微分、复合相容的提升。

\ProseParagraphBreak
说一个三角范畴具有 DG 描述，还必须解释它的移位与好三角来自哪里。为此可以用 DG 模表达闭性条件。一个右 DG \(\mathcal E\)-模是 DG 函子 \(\mathcal E^{\mathrm{op}}\to\mathrm{Ch}_{\mathrm{dg}}(k)\)，其中反范畴的复合采用上一小节的符号。

\ProseParagraphBreak
先把这个定义翻译到只有一个对象的情形，此时 \(\mathcal E\) 就是 DG 代数 \(A\)。通常的右 DG \(A\)-模写作右作用 \(ma\)，满足
\[
 \mathrm d(ma)=(\mathrm dm)a+(-1)^{|m|}m\,\mathrm da.
\]
它对应的 DG 函子在齐次元素上采用
\[
 M(a^{\mathrm{op}})(m)=(-1)^{|a||m|}ma.
\]
这个符号与 \(a^{\mathrm{op}}b^{\mathrm{op}}=(-1)^{|a||b|}(ba)^{\mathrm{op}}\) 相配。例如 \(|a|=|b|=|m|=1\) 时，连续作用给出 \(M(a^{\mathrm{op}})M(b^{\mathrm{op}})(m)=-mba\)，与反代数中乘积的作用相同。因而右作用的通常写法与反范畴函子的写法之间，需要这一次符号转换。

\ProseParagraphBreak
对两个右 DG \(\mathcal E\)-模 \(M,N\)，次数 \(r\) 的态射是一族次数 \(r\) 的 \(k\)-线性映射 \(\eta_Y:M(Y)\to N(Y)\)，满足 \(N(f)\eta_Y=(-1)^{r|f|}\eta_{Y'}M(f)\)，其中 \(f:Y'\to Y\) 为 \(\mathcal E\) 中的任意齐次态射。在单对象情形，把上面的转换式代入，两边的次数符号抵消，便得到通常的右线性条件 \(\eta(ma)=\eta(m)a\)。态射复形的微分仍为\eqref{b4:eq:dg-hom-differential}；模的移位、映射锥逐对象构造。对齐次态射 \(f\)，移位作用为 \(M[1](f)=(-1)^{|f|}M(f)\)；对闭的零次模态射 \(u:M\to N\)，锥上的作用为 \(\operatorname{diag}(N(f),(-1)^{|f|}M(f))\)。这些符号使作用与移位及锥的微分相容。

\ProseParagraphBreak
可表模 \(h_X\) 把 \(Y\) 送到 \(\underline{\operatorname{Hom}}_{\mathcal E}(Y,X)\)，齐次 \(f:Y'\to Y\) 对 \(a:Y\to X\) 的作用为 \((-1)^{|f|\cdot|a|}af\)。在单对象情形，\(h_*=A\)；若 \(a,f\) 都为一次，函子作用是 \(h_*(f^{\mathrm{op}})(a)=-af\)。再按上面的转换式还原通常右作用，得到的正是右正则模的乘法 \(af\)。这说明 Yoneda 公式中的负号来自反范畴约定，而不是把右乘法本身改成负的乘法。

\begin{definition}\label{b4:def:dg-enhancement}
设 \(\mathcal E\) 为交换环 \(k\) 上的 DG 范畴。若其可表右 DG 模及与它们同构的对象在 DG 模的同伦范畴中组成一个包含零对象、对正负移位及零次闭态射的映射锥封闭的满子范畴，称 \(\mathcal E\) 为\term[yu-sanjiao-DGfanchou]{预三角 DG 范畴}[pretriangulated DG category]。设 \(\mathcal T\) 为 \(k\)-线性三角范畴；一对数据 \((\mathcal E,\Phi)\)，其中 \(\mathcal E\) 预三角而 \(\Phi:H^0(\mathcal E)\to\mathcal T\) 为三角等价，称为 \(\mathcal T\) 的\term[DGzengqiang]{DG 增强}[DG enhancement]。
\end{definition}

先在全部右 DG 模的同伦范畴中构造三角结构。移位与锥逐对象形成复形；上述作用符号保证它们仍是 DG 模。\cref{b4:thm:homotopy-triangulated}证明中的锥映射及同伦矩阵也与模作用相容，因此同一验证给出这里的三角公理。DG Yoneda 函子则在态射复形上给出同构
\[
\underline{\operatorname{Hom}}_{\mathrm{DGMod}(\mathcal E)}(h_X,h_Y)
\cong\underline{\operatorname{Hom}}_{\mathcal E}(X,Y).
\]
其中 \(\mathrm{DGMod}(\mathcal E)\) 表示右 DG 模的 DG 范畴；该同构把模态射送到它在 \(1_X\) 处的值，逆映射把齐次 \(g:X\to Y\) 送到逐分量的后复合 \(a\mapsto ga\)。带符号的自然性与 Leibniz 法则保证这两个对应与微分相容。取零次上同调，便得到 \(H^0(\mathcal E)\) 到 DG 模同伦范畴的全忠实函子。

\ProseParagraphBreak
若 \(\mathcal E\) 预三角，这个函子的本质像包含零对象，并对正负移位及锥封闭。它也对有限双积封闭，因为零态射 \(h_X[-1]\to h_Y\) 的锥就是 \(h_Y\oplus h_X\)。把 DG 模同伦范畴的好三角限制到该满子范畴，便得到三角结构：三角的补全、旋转及八面体中所需的对象都是这些移位和锥，因而仍在本质像中，三角态射的各分量则由满性保留。DG Yoneda 的全忠实性给出 \(H^0(\mathcal E)\) 与这个本质像的等价，由此将三角结构传回 \(H^0(\mathcal E)\)。这才是定义中 DG 增强所用的三角结构。DG 模的这些构造也可查 Keller 的《Deriving DG categories》\cite[Sections 1.2 and 2.2]{Keller1994DerivingDG}。一般 DG 范畴可能既没有移位对象，也没有锥对象，不能仅凭 DG 模同伦范畴具有三角结构，就断言其 \(H^0\) 是三角范畴。

\begin{example}\label{b4:exm:bounded-projective-enhancement}
设 \(R\) 为含幺环。取所有有界、每项有限生成投射的左 \(R\)-模复形以及它们的整个 Hom 复形，得到 \(\mathrm{Ch}^b_{\mathrm{dg}}(R\text{-}\mathrm{proj})\)。移位仍在其中；映射锥的项为两个有限生成投射模的直和，也仍在其中。因此它是预三角的，\(H^0\) 是 \(K^b(R\text{-}\mathrm{proj})\)。下一节把这个范畴识别为完美复形的导出范畴模型。这是一个直接可写出的增强，而非从抽象三角公理中反求出来的结构。
\end{example}
